\section{Architecture and Design}
\label{sec:arch}

\subsection{Design Principles}
Four principles follow from Section~\ref{sec:math}. (1)~\emph{The gate is
derived, not tuned}: the controller compares a calibrated error estimate with
a cost ratio (Theorem~\ref{thm:gate}). (2)~\emph{Acceptance needs evidence}:
no answer is marked verified unless an exact verifier accepted it; consensus
and judge scores are assurance levels, not proofs
(Theorems~\ref{thm:precision} and \ref{thm:plurality}). (3)~\emph{Every loop is
bounded}: $N$, $R$ and the thinking budget are configuration limits enforced
by the gateway (Lemma~\ref{lem:termination}). (4)~\emph{Self-knowledge is
audited}: the self-model is built from traces plus a random audit stream and
reports unmeasured capabilities as such (Proposition~\ref{prop:audit}).

\subsection{Concept-to-Component Mapping}
Fig.~\ref{fig:mapping} gives the mapping. The orchestrator (\texttt{nafs-core})
plays al-Ghazali's king: it owns the request and consults the controller. The
controller (\Gate) plays the vizier: it arbitrates between the fast draft and
deliberation on a cost--risk basis. The fast path (\Sone) is the impulsive
draft; the slow path (\Stwo) is deliberation; the \Hub\ supplies external
evidence; the \Ledger\ performs the nightly accounting; \Distil\ is the long-run
change of disposition. The gate's cheap features play the role of Avicenna's
estimative faculty, a fast particular judgement that is useful and
error-prone, which is exactly why it is calibrated and audited.

\begin{figure}[t]
\centering
% Concept-to-component mapping (analogy only); width ~8.6cm
\begin{tikzpicture}[x=1cm,y=1cm]
\foreach \i/\txt/\col/\eng in {
 0/{al-nafs al-ammarah\\(Q 12:53) impulsive self}/navy/{\Sone\ draft: fast, unverified},
 1/{al-nafs al-lawwamah\\(Q 75:2) self-reproaching self}/crimson/{\Gate\ flag or verifier rejection},
 2/{al-nafs al-mutma'innah\\(Q 89:27) self at rest}/leaf/{Accepted answer with assurance level},
 3/{'aql (al-Ghazali: vizier)\\deliberating intellect}/amber/{\Gate: cost--risk arbitration; \Stwo\ deliberation},
 4/{qalb (al-Ghazali: king)\\integrating executive}/violet/{Orchestrator \texttt{nafs-core}: owns the loop},
 5/{wahm, dhakira (Ibn Sina)\\estimation, memory}/teal/{Gate features; \Mem},
 6/{muhasaba, muraqaba\\self-accounting, vigilance}/teal/{\Ledger: audits, calibration},
 7/{tazkiyat al-nafs\\purification}/slate/{\Distil: verified traces $\to$ LoRA}}
{ \node[sblk=\col,minimum width=2.9cm,text width=2.75cm] (c\i) at (0,-0.78*\i) {\txt};
  \node[sblk=\col,minimum width=3.9cm,text width=3.75cm] (e\i) at (4.75,-0.78*\i) {\eng};
  \draw[arr] (c\i)--(e\i); }
\node[font=\tiny\bfseries,text=slate] at (0,0.62) {Classical concept (source)};
\node[font=\tiny\bfseries,text=slate] at (4.75,0.62) {Engineering component (analogy)};
\end{tikzpicture}

\caption{Concept-to-component mapping. Every arrow is an analogy; component
names in code are the engineering names on the right.}
\label{fig:mapping}
\end{figure}

\subsection{Request Lifecycle}
Fig.~\ref{fig:flow} shows one request and Fig.~\ref{fig:nafs-automaton} its
state automaton. The \Mem\ assembles the prompt from the conversation and any
retrieved memory with provenance. \Sone\ produces a draft in non-thinking mode
together with $k$ cheap samples in one request (the \texttt{n} parameter of
the OpenAI-compatible API, so the prompt is prefilled once). The \Gate\
computes features: mean and minimum token log-probability of the draft,
agreement among the $k$ samples, the domain class, and, for verifiable tasks,
the verifier outcome on the draft. A calibrated model maps the features to
$\hat e_1$; the gate escalates if $L(\hat e_1-\hat e_2)>\lambda\Delta c$. A
verifiable task whose draft fails verification is never eligible for the
keep branch (this closes the draft's defect~10, where the router pseudo-code
disagreed with its prose). On escalation, \Stwo\ runs draft-blind in thinking
mode with budget $b$ and sequential sampling: each candidate is verified and
the first accepted one is returned; a rejected candidate's evidence (failing
test, counterexample, schema error) is fed into the next round, up to $R$
rounds. If the budget is exhausted the reply is an explicit abstention with
the best candidate and an uncertainty notice. Every reply carries an
\emph{assurance level}: \texttt{verified} (exact verifier passed),
\texttt{consensus} (soft agreement only), \texttt{unverified} (kept fast
draft on a non-verifiable task) or \texttt{abstained}. The decision record of
the gate and the trace of the loop are written for the \Ledger.

\begin{figure}[t]
\centering
% Request lifecycle: fast path, gate, deliberation, verification; width ~8.6cm
\begin{tikzpicture}[x=1cm,y=1cm]
\node[blk=navy,text width=1.1cm]   (in)  at (0.1,0)    {Request $x$};
\node[blk=teal,text width=1.45cm]   (ctx) at (1.95,0)  {\Mem\\history, cache};
\node[blk=navy,text width=1.45cm]   (s1)  at (3.95,0)  {\Sone\\draft $+$ $k$ samples};
\node[dm,minimum width=1.6cm]      (g)   at (6.2,0)  {\Gate\\escalate?};
\node[blk=violet,text width=1.45cm] (s2)  at (6.2,-1.7) {\Stwo\\draft-blind, budget $b$};
\node[blk=amber,text width=1.45cm]  (v)   at (3.95,-1.7) {\Hub\\$(r,\varphi)$};
\node[blk=slate,text width=1.45cm]  (tr)  at (1.95,-1.7) {Trace, decision record};
\node[blk=leaf,text width=1.45cm]   (out) at (6.2,-3.3) {Reply $+$ assurance level};
\draw[arr] (in)--(ctx); \draw[arr] (ctx)--(s1); \draw[arr] (s1)--(g);
\draw[arr] (g.south)--node[lab,right]{escalate}(s2.north);
\draw[arr] (g.east)--++(0.45,0)|-node[lab,pos=.25,right]{keep}(out.east);
\draw[arr] (s2)--node[lab,above=1pt]{candidate}(v);
\draw[dsh] (v.south east) to[bend right=28] node[lab,below]{rejected: retry $\le R$} (s2.south west);
\draw[arr] (v.south)--++(0,-0.75)-|node[lab,pos=.25,below]{accepted}(out.west);
\draw[dsh] (s1.south)--node[lab,right]{verify kept draft}(v.north);
\draw[arr] (v)--(tr);
\end{tikzpicture}

\caption{Request lifecycle. Solid arrows are the main path; dashed arrows are
the verification and retry edges.}
\label{fig:flow}
\end{figure}

\begin{figure}[t]
\centering
% Trace-state automaton of one request (A/L/M/T/X); width ~8.4cm
\begin{tikzpicture}[x=0.93cm,y=1cm,>=Stealth]
\node[st=navy]    (A) at (0,0)   {$\mathsf A$};
\node[st=crimson] (L) at (2.6,0) {$\mathsf L$};
\node[st=violet]  (M) at (5.2,0) {$\mathsf M$};
\node[st=leaf]    (T) at (2.6,-1.9) {$\mathsf T$};
\node[st=slate]   (X) at (7.6,-1.9) {$\mathsf X$};
\draw[arr] (-1.1,0) -- (A) node[lab,pos=0.3,above]{request};
\draw[arr] (A) -- node[lab,above]{flagged} (L);
\draw[arr] (L) -- node[lab,above]{retry} (M);
\draw[arr] (A) -- node[lab,left,pos=.55]{kept $\wedge$ verified} (T);
\draw[arr] (M) -- node[lab,below right,pos=.35]{accepted} (T);
\draw[arr] (M) to[bend left=35] node[lab,below]{rejected, $t<R$} (L);
\draw[arr] (M) -- node[lab,above right,pos=.4]{budget out} (X);
\node[font=\tiny,align=left,anchor=north west,text width=7.9cm] at (-1.05,-2.6)
 {$\mathsf A$ draft (fast path); $\mathsf L$ flagged by gate or verifier; $\mathsf M$ deliberate retry conditioned on verifier evidence;
  $\mathsf T$ accepted (assurance level attached); $\mathsf X$ abstain with uncertainty notice.};
\end{tikzpicture}

\caption{Trace-state automaton of one request. The three Qur'anic stage names
are used as mnemonics for $\mathsf A$, $\mathsf L$ and $\mathsf T$; the
automaton is a bounded retry loop (Lemma~\ref{lem:termination}).}
\label{fig:nafs-automaton}
\end{figure}

\subsection{Components}
\label{sec:components}
\begin{itemize}
\item \textbf{\Gw\ (\texttt{nafs-gateway}).} HTTP API (OpenAPI 3.1), SSE
streaming of tokens and of controller events, WebSocket chat rooms,
authentication, rate limits, Prometheus metrics.
\item \textbf{Orchestrator and \Gate\ (\texttt{nafs-core}).} Implements the
automaton, the gate, sequential verified sampling, budget forcing, timeouts
and fallbacks; emits the decision record.
\item \textbf{\Sone\ and \Stwo\ backends.} Two OpenAI-compatible model
servers (or one server with two modes in the single-model profile), reached
over HTTP with retry and backpressure.
\item \textbf{\Hub\ (\texttt{nafs-verify}).} Exact verifiers: code sandbox
(subprocess with CPU, memory, wall-clock and network limits), numeric and
symbolic mathematics (reference-answer equality, SymPy via a sidecar), JSON
Schema validation, format checks. Soft verifiers: rubric judge (the \Stwo\
model in non-thinking mode at temperature~0) and sample consensus. Each
verifier publishes $(r,\varphi)$ from mutation tests.
\item \textbf{\Mem\ (\texttt{nafs-memory}).} Conversation store (SQLite),
prompt assembly within a context cap, prompt-cache affinity hints for the
model server, episodic trace store (JSONL), and a retrieval index built with
Qwen3-Embedding-0.6B when retrieval is enabled.
\item \textbf{\Ledger.} Offline job that aggregates traces and the $5\%$ audit
stream into per-domain calibration tables (escalation rate, $\hat e_1$
calibration error, verifier acceptance rates, latency quantiles), keyed by
model and adapter hash; feeds the dashboard and the gate's recalibration.
\item \textbf{\Harness\ (\texttt{nafs-eval}) and performance console
(\texttt{nafs-bench}).} Dataset loaders, metric definitions, paired
statistical tests, load generation, and report export.
\item \textbf{Web UI (\texttt{nafs-web}).} Chat room with the assurance level
and the reasoning mode visible per message, a performance-test console, and an
evaluation dashboard; server-rendered HTML with htmx for streaming updates.
\item \textbf{\Distil\ (Python, GPU).} Dataset builder from verified traces,
Unsloth LoRA training, promotion gate, GGUF export.
\end{itemize}

\begin{figure}[t]
\centering
% System architecture (single-node CPU baseline); width ~8.6cm
\begin{tikzpicture}[x=1cm,y=1cm]
\node[blk=slate,text width=1.4cm] (ui) at (0.1,0) {Web UI\\chat, bench, eval};
\node[blk=slate,text width=1.4cm] (cli) at (0.1,-1.2) {API clients\\SDK, harness};
\node[blk=navy,text width=1.8cm] (api) at (3.1,0) {\texttt{nafs-gateway}\\HTTP, SSE, WS};
\node[blk=navy,text width=1.8cm] (core) at (3.1,-1.2) {\texttt{nafs-core}\\\Gate, loop};
\node[blk=amber,text width=1.8cm] (hub) at (3.1,-2.4) {\texttt{nafs-verify}\\sandbox, math, schema};
\node[blk=teal,text width=1.8cm] (mem) at (3.1,-3.6) {\texttt{nafs-memory}\\SQLite, traces};
\node[blk=violet,text width=1.8cm] (b1) at (6.35,0) {\Sone\ server\\OpenAI-compatible};
\node[blk=violet,text width=1.8cm] (b2) at (6.35,-1.2) {\Stwo\ server\\thinking mode};
\node[blk=crimson,text width=1.8cm] (sb) at (6.35,-2.4) {Code sandbox\\subprocess, limits};
\node[blk=leaf,text width=1.8cm] (obs) at (6.35,-3.6) {Metrics, logs\\Prometheus, JSONL};
\begin{scope}[on background layer]
\node[grp=navy,fit=(api)(core)(hub)(mem),inner sep=6pt] (gw) {};
\node[grp=violet,fit=(b1)(b2),inner sep=6pt] (srv) {};
\end{scope}
\node[grplab=navy,anchor=south west] at (gw.north west) {Rust workspace};
\node[grplab=violet,anchor=south west] at (srv.north west) {vLLM-CPU or llama-server};
\draw[arr] (ui)--(api); \draw[arr] (cli)--(api);
\draw[arr] (api)--(core); \draw[arr] (core)--(hub); \draw[arr] (core)--(mem);
\draw[arr] (core.east)--node[lab,above,pos=.45]{HTTP}(b1.west);
\draw[arr] (core.east)--(b2.west);
\draw[arr] (hub)--(sb);
\draw[dsh] (mem)--(obs);
\end{tikzpicture}

\caption{System architecture, single-node CPU baseline. The GPU scale-out
variant replaces the model servers with vLLM on CUDA and keeps every other
component unchanged.}
\label{fig:arch}
\end{figure}

\subsection{Model Profiles}
\label{sec:models}
Table~\ref{tab:profiles} defines three profiles. Profile~A is the reference:
two dense Qwen3 checkpoints, one per role, chosen because dense
grouped-query attention has the most mature CPU kernels, official GGUF files
exist, and both are on vLLM's validated CPU list. Profile~B serves one
switchable Qwen3-4B for both roles, halving resident memory at the price of a
slower fast path; the gate then decides \emph{whether to think}, exactly the
decision studied by the adaptive-thinking literature. Profile~C is the
upgrade path to Qwen3.5, which is stronger and multimodal but uses hybrid
linear attention whose CPU kernels are newer and, on AVX2 without
AVX-512-BF16 or AMX, unaccelerated in vLLM; Phase~0 of the roadmap benchmarks
it against Profile~A on the target machine before any commitment. The
upgrade decision is a go/no-go gate with a measurable criterion.

\begin{table}[t]
\caption{Model profiles (Hugging Face IDs verified 2026-10-18; sizes are GGUF file sizes)}
\label{tab:profiles}
\centering\scriptsize
\begin{tabular}{@{}p{0.7cm}>{\raggedright\arraybackslash}p{3.6cm}>{\raggedright\arraybackslash}p{3.6cm}@{}}
\toprule
Profile & \Sone & \Stwo \\
\midrule
A (ref.) & \texttt{Qwen/\allowbreak Qwen3-1.7B}, GGUF Q4\_K\_M 1.3\,GB (\texttt{unsloth/\allowbreak Qwen3-1.7B-GGUF}) or Q8\_0 1.8\,GB (\texttt{Qwen/\allowbreak Qwen3-1.7B-GGUF}); thinking off & \texttt{Qwen/\allowbreak Qwen3-4B-\allowbreak Thinking-2507}, GGUF Q4\_K\_M 2.5\,GB (\texttt{unsloth/\allowbreak Qwen3-4B-\allowbreak Thinking-2507-\allowbreak GGUF}); 262K context \\
B (min.) & \texttt{Qwen/\allowbreak Qwen3-4B} Q4\_K\_M 2.5\,GB (\texttt{Qwen/\allowbreak Qwen3-4B-GGUF}), \texttt{enable\_thinking=false} & same weights, \texttt{enable\_thinking=true} \\
C (upgr.) & \texttt{Qwen/\allowbreak Qwen3.5-2B} Q4\_K\_M 1.3\,GB (\texttt{unsloth/\allowbreak Qwen3.5-2B-GGUF}) & \texttt{Qwen/\allowbreak Qwen3.5-4B} Q4\_K\_M 2.7\,GB (\texttt{unsloth/\allowbreak Qwen3.5-4B-GGUF}) \\
\midrule
\multicolumn{3}{@{}p{8.2cm}@{}}{Judge: the \Stwo\ model in non-thinking mode. Embeddings: \texttt{Qwen/\allowbreak Qwen3-Embedding-0.6B-GGUF} Q8\_0. All Apache-2.0. GPU scale-out: the same IDs in BF16 or AWQ under vLLM CUDA, with \texttt{Qwen/Qwen3-8B} as the \Stwo\ option.} \\
\bottomrule
\end{tabular}
\end{table}

\subsection{Deployment Topology}
The CPU baseline is a single node: one gateway process, one or two model
servers pinned to the physical cores (the gateway reserves none, since the
gateway's own CPU use is small and bursty), a sandbox runner, SQLite and
JSONL on local disk. The model servers are either \texttt{llama-server}
(macOS or any x86) or vLLM-CPU (Linux, AVX-512 recommended), both behind the
same OpenAI-compatible API, so the Rust code has one backend trait and two
thin implementations. The optional GPU scale-out variant moves the model
servers to vLLM on CUDA with both models resident, raises $N$ and the
thinking budget, and enables EAGLE or multi-token-prediction speculative
decoding; the gateway, verifiers, memory and UI are unchanged.

\subsection{What Changed from the Draft}
The draft's twelve validation results are retained where exact (precision,
coverage, cascade cost, router threshold, consensus, value-score algebra,
coupling, audit, measurement floors, distillation noise, capacity, cascade
identity) and are restated here with explicit assumptions. We added the
miscalibration bound, the frontier proposition, the optimal-stopping cap, the
two refinement-loop theorems, the plurality limit, the roofline and the
thinking-budget condition. We removed the mapping of the Active Intellect to
``external verifiers'' and of \emph{ilham} to any component, following the
guidance in Section~\ref{sec:islamic}; we replaced the GPU roster with CPU
profiles; and we turned the draft's ten defects into requirements with tests
in the specification.
